% !TEX root = ../main.tex
\section{Introduction}
\label{sec:introduction}

Many standard econometric models assume that aggregate behavior mirrors individual behavior (up to a random error term). 
For instance, a traditional wage equation might specify that an individual's expected wage increases by a population-constant parameter $\beta$ for each additional year of schooling, assuming homogeneous returns to education across the population. 
However, this assumption of parameter homogeneity is often overly restrictive. 
Consumer preferences, for example, exhibit significant variation; while a specific vehicle attribute (such as a distinct color) might increase the purchase probability for some consumers, it could deter others. 
Consequently, accurately modeling such unobserved preference heterogeneity is crucial.

Random coefficient models provide a robust framework for incorporating this heterogeneity. 
These models have become foundational across various disciplines, including labor economics, industrial organization, and asset pricing \citep[e.g.][]{BerryLevinsohnPakes1995, CarneiroHeckmanVytlacil2011, BelzilHansen2007, KoijenYogo2019}. 
For a comprehensive overview of random coefficient models in the context of panel data, see \citet{HsiaoPesaran2008}. 
Another prominent class within the random coefficient framework consists of random coefficient discrete choice models, which extend traditional discrete choice frameworks like the probit, logit, and multinomial logit models \citep[see][]{CameronTrivedi2005}. 
A standard example is the mixed logit model \citep[chap.~6]{Train2009}.

The mixed logit generalizes the standard multinomial logit by allowing the marginal utility coefficients to vary randomly across individuals. 
A major theoretical advantage of the mixed logit over the multinomial logit is that it relaxes the restrictive independence of irrelevant alternatives (IIA) assumption \citep{Train2009}. 
As a result, the mixed logit has seen widespread empirical application \citep[e.g.][]{NicoletNegenbornAtasoy2022,YinCherchi2024,ZhouWangYangGanLiuZi2025}. 
When the mixing distribution of these random coefficients is assumed to be discrete, the mixed logit simplifies to a latent class multinomial logit model \citep{Train2009}, as utilized by \citet{WangYangSongSia2014} and \citet{LiGaoYangLiu2022}. 

The literature proposes numerous estimation procedures for this model. 
\citet{Train2009} outlines simulated maximum likelihood, the expectation-maximization (EM) algorithm \citep[further detailed in][]{Train2008}, and hierarchical Bayes methods using Markov chain Monte Carlo (MCMC) \citep[see also][]{BurdaHardingHausman2008,Rossi2012}. 
Other approaches include the simulated method of moments \citep{McFadden1989} and minimum distance estimators \citep{BeranMillar1994}. 
However, as highlighted by \citet{FoxKimRyanBajari2011}, many of these traditional methods suffer from computational burdens and convergence issues. 
To address this, after a preliminary paper by \citet{BajariFoxRyan2007}, \citet{FoxKimRyanBajari2011} proposed a computationally tractable estimator, commonly referred to as the FKRB estimator.
Their estimator has gained significant traction in the recent literature \citep[see for example][]{HoudeMyers2019,BlundellGowrisankaranLanger2020}.

The fundamental insight of the FKRB approach is that the mixed logit can be recast such that the parameters enter the model linearly. 
Consequently, these parameters can be estimated using ordinary least squares (OLS).
However, the parameters have to satisfy some constraints. 
Building on this, \citet{HeissHetzeneckerOsterhaus2022} demonstrated that the constraints inherent to the FKRB framework render it mathematically equivalent to a nonnegative LASSO (NNL) estimator. 
Leveraging this equivalence, they further extended the approach to a nonnegative elastic net estimator.

\citet{Train2016} proposed an alternative refinement to the FKRB estimator by approximating the distribution of random coefficients using splines, polynomials, or step functions, rather than relying on a discrete mixture. 
Researchers frequently employ discrete mixture estimators primarily for their computational tractability, even when the true underlying distribution is believed to be continuous. 
As shown by \citet{McFaddenTrain2000}, a discrete distribution with a sufficiently dense support can approximate any well-behaved continuous distribution to an arbitrary degree of accuracy.

The FKRB estimator can be seen in both a parametric and a nonparametric framework.
The nonparametric FKRB estimator is a sieve estimator \citep{HeissHetzeneckerOsterhaus2022, FoxKimYang2016}, specifically a (penalized) sieve minimum distance ((P)SMD) estimator.
The SMD estimator has been described by \citet{AiChen2003,NeweyPowell2003}.
\citet{ChenPouzo2012} introduced the PSMD estimator, which is a generalization of the SMD estimator.
Sieves were first introduced by \citet{Grenander1981}.
Many papers have used some form of sieve estimation, including \citet{GallantNychka1987}, who proposed an estimator that can be written as a sieve maximum likelihood estimator; \citet{WhiteWooldridge1991} and \citet{ChenShen1998}, who developed theory on sieve estimation in the context of dependent data; and \citet{Andrews1991} and \citet{Newey1997}, who developed much of the asymptotic theory on series estimators.
An extensive overview of sieve estimation is given by \citet{Chen2007}.

Beyond estimation, conducting valid statistical inference remains a primary concern. 
\citet{FoxKimRyanBajari2011} explore several methods for constructing confidence intervals for their estimator. 
They initially suggest utilizing clustered heteroskedasticity-robust standard errors derived from an unconstrained OLS procedure, clipping the resulting intervals to ensure they remain within the $[0,1]$ bound. 
While this guarantees correct asymptotic coverage, the resulting intervals have been observed to be somewhat conservative, covering the true parameter more frequently than strictly necessary, which results in too many type II errors. 
They also evaluate resampling techniques, such as the standard bootstrap \citep{Efron1979} and the $m$-out-of-$n$ bootstrap, but find limited evidence that these yield accurate coverage. 
Finally, they try Tikhonov regularization, noting a marginal improvement in inference quality. 
This essentially extends the estimator to that of \citet{HeissHetzeneckerOsterhaus2022}.
However, \citet{HeissHetzeneckerOsterhaus2022} do not provide an inference procedure.

\citet{FoxKimRyanBajari2011} indicate that an explanation for the conservativeness of their inference procedures may lie in the fact that the true parameters are often close to the boundary of the parameter space.
On this boundary, asymptotic normality breaks down \citep{Andrews1999}.
\citet{Andrews1999,Andrews2001,Andrews2002} give the asymptotic distribution for related cases with a true parameter on the boundary.
\citet{AndrewsGuggenberger2010Applications} propose a subsampling procedure for another related case.
According to \citet{FoxKimRyanBajari2011}, these results however are unfortunately not sufficient to cover the FKRB estimator.
\citet{Li2025The,Li2025Inference} develops infernce for constrained extremum estimators and gives the FKRB estimator as an example.

\citet{ChenPouzo2015} established theoretical results for inference on a very broad class of PSMD estimators.
A downside of their results is that they are exclusive to linear parameter spaces, which do not include bounded spaces.
Therefore, the answer to the boundary problem may not lie in \citet{ChenPouzo2015}.

Although we acknowledge that the FKRB estimator does not belong to the class of \citet{ChenPouzo2015}, we note that it is conjectured by \citet{FoxKimYang2016} that similar results to those of \citet{ChenLintonVanKeilegom2003} hold for the FKRB estimator.
Therefore, this thesis investigates the performance of the sieve Wald and quasi-likelihood ratio (QLR) statistics from \citet{ChenPouzo2015} for conducting inference on the FKRB estimator.
We will compare to the inference procedures originally proposed by \citet{FoxKimRyanBajari2011}.

The thesis studies the coverage of confidence intervals stemming from the aforementioned procedures using a Monte Carlo simulation.
We consider seven different data generating processes (DGPs).
These DGPs differ merely in the distribution of random coefficients.
The different distributions are chosen to investigate multiple levels of boundary problems.
We will examine how the inference procedures hold up in two scenarios.
In one scenario we estimate the whole distribution, in the other we only estimate the mean of the distribution.

At small grid sizes the simulation reveals that the standard FKRB inference procedure is often statistically indistinguishable from an optimal procedure for off-boundary distributions, but conservative for on-boundary distributions.
At large grid sizes, however, the procedure is incorrect for off-boundary distributions, in the sense that the empirical coverage is lower than the nominal; for on-boundary distributions the coverage is correct.
The sieve Wald and QLR procedures are on the conservative side for both small and big grids and both off- and on-boundary distributions.

The thesis will proceed as follows. 
\Cref{sec:theory} lays down the theoretical framework in which the FKRB estimator resides.
\Cref{sec:mxl} explains the mixed logit model, \cref{sec:FKRBestimator} derives the FKRB estimator, alongside some of its extensions and equivalent forms, and \cref{sec:Inference} formulates the inference procedures that are researched in this thesis.
\Cref{sec:simulation} describes the Monte Carlo simulation procedure.
\Cref{sec:dgp} lays out the DGPs, including the different random coefficient distributions, and \cref{sec:analysis} specifies the techniques we use to analyze the results.
\Cref{sec:results} gathers all the results and \cref{sec:conclusion} concludes the thesis.
